\section{Exact arithmetic and the finite formal certificate}
\label{sec:verification}

The fixed finite numerical input consists of
the matrix~\eqref{eq:M}, its six rank-one row and column matrices,
and the six rational values in~\eqref{eq:row-determinants} and
\eqref{eq:column-determinants}. All entries in the tensor and both
decompositions are integers. The argument in Section~\ref{sec:witness}
proves the tangent ranks from these determinants without floating-point
rank decisions.

\subsection{Independent exact arithmetic checks}

Two separate exact implementations verify the explicit data. One reads
the two lists of integer factors; the other constructs them directly
from~\eqref{eq:M}. They verify the equality of the resulting tensors,
nonzero and distinct summands, distinct decompositions, and the Jacobian
ranks $45$ and $80$. The first additionally checks pairwise distinct
projective directions, all three pair-product ranks, and the rank-five
third flattening explicitly.
The coordinate Jacobian in the smaller format has $45$ rows and
$5(3+3+5)=55$ columns; after padding it has $216$ rows and
$5(6+6+6)=90$ columns.

Exact elimination is performed modulo the primes $1{,}000{,}003$ and
$101$, respectively. A nonzero minor modulo a prime is a nonzero
integer minor, so these computations give rigorous lower bounds for
the ranks over $\Q$ and $\C$. The prime used by the first implementation
is itself verified by integer trial division. For the matching upper
bounds, each nonzero pure summand has two independent scaling
directions in the factor-coordinate kernel, for example
\[
 (\dot a_i,\dot b_i,\dot c_i)=(a_i,-b_i,0),
 \qquad (a_i,0,-c_i).
\]
Across five summands these give ten independent kernel vectors.
The bounds are therefore $55-10=45$ and $90-10=80$, matching the
modular lower bounds.

The verification includes altered-data tests that break tensor equality,
pair-product independence, or distinctness of the decompositions.
These are required to be rejected. The positive checks and the
altered-data checks have identical prescribed behavior with and without
Python optimization, and with two choices of hash seed. All decisive
arithmetic uses integers, rational numbers, or finite prime fields.

\subsection{The Lean statement and its limits}

The accompanying Lean development proves a single finite certificate
whose assertions are the following:
\begin{enumerate}
 \item The two integer factor lists have the same sum, both in format
 $3\times3\times5$ and after padding to $6\times6\times6$.
 \item In each list the summands are nonzero and pairwise
 nonproportional, and the two lists give different sets of directions;
 these assertions also hold after padding.
 \item The two third-mode flattening matrices, before and after
 padding, have rank five over $\C$.
 \item The six determinants are exactly the rational numbers in
 \eqref{eq:row-determinants} and~\eqref{eq:column-determinants}.
 \item Both coordinate Jacobians have rank $45$ over $\C$ in the
 smaller format and rank $80$ over $\C$ after padding.
\end{enumerate}
The formal rank lower bounds use integer right-inverse identities
with scale factor $36$. The ten explicit scaling directions supply
the upper bound for the padded Jacobian. The formalization uses Lean
and Mathlib, with the certificate fixed by the displayed integer data.

These are statements about concrete arrays, determinants and ranks.
The formal theorem does not state a tensor border-rank bound, identify
the geometric tangent space, prove the dimension of a secant variety,
or establish any fact about normalization, conductors, $S_2$,
Cohen--Macaulayness, or vertex depth. The passage from the finite
certificate to Theorem~\ref{thm:main}, including the uniform extension
to all ranks and formats in~\eqref{eq:main-range}, is the written
proof in the preceding sections. In particular,
neither the tangent ranks alone nor nonnormality alone imply the
Cohen--Macaulay conclusion.


\subsection{The uniform construction}

The extension to arbitrary $r$ uses the identities in
Lemma~\ref{lemma:avoidance} and the block comparisons in
Lemma~\ref{lemma:extended-tangencies}. Identity
\eqref{eq:avoidance-determinant} can also be verified coefficient by
coefficient in $\Q[t]$; its nonvanishing for every integer $t\geq4$
is an immediate consequence of the displayed factorization.
Thus no search for witnesses or recognition of numerical data is
needed as $r$ grows.

For additional finite corroboration, exact arithmetic checks the
$15$ pairs $(r,p)$ with $5\leq r\leq9$ and $3\leq p\leq r-2$,
putting $q=r+1-p$. Each seed in format $p\times q\times r$
and each padded example in format $p\times(q+1)\times r$ is checked,
as is the extension from $4\times5\times8$ to $5\times5\times8$.
These are $31$ pairs of decompositions. Equality and distinctness
are tested over the integers, while pair-product and Jacobian ranks
are checked modulo $101$, with the displayed scaling kernels giving
matching characteristic-zero upper bounds. Altering equality,
distinctness, independence, or an avoidance vector is required to
cause rejection, with the same behavior under Python optimization.
These checks concern exactly the listed instances. The all-$r$
statement, the smoothness of the incidence charts, the conductor
argument, and the local depth conclusions rest on their written proofs.
